% =========================
% 2026 ICM Problem F Report Template (≤25 pages)
% Page 1 MUST be Summary Sheet
% Header on EVERY page: Team #<CONTROL>, Page i of 25
% =========================

\section*{Summary Sheet (Page 1)}
\textbf{Title:} \\
\textbf{Careers Chosen (STEM/Trade/Arts):} \\
\textbf{Institutions \& Programs:} \\
\textbf{One-sentence problem framing:} \\
\textbf{Modeling approach (1–2 lines):} \\
\textbf{Key quantitative findings (3 bullets):} \\
\textbf{Recommendations (3 institutions, 1–2 bullets each):} \\
\textbf{Validation \& sensitivity (1 line):} \\
\textbf{Keywords:} (3–6)

\newpage
\tableofcontents
\newpage

\section{Introduction}
\subsection{Background and Motivation}
% ≤8 sentences total background across the whole report.
\subsection{Problem Restatement (Task Mapping)}
% Map explicitly to the prompt:
% (i) 3 careers, (ii) data-informed future model, (iii) 3 institutions/programs & recommendations,
% (iv) optional: other success metrics & generalization.
\subsection{Contributions and Roadmap (Our Work)}
% Insert Fig.1 workflow overview and 4–6 bullet contributions.

\section{Assumptions and Justifications}
% Table format; each assumption: statement + 1-line justification + impact note.

\section{Notations}
% One table for symbols, units, and meaning (define once).

\section{Data and Evidence Base}
\subsection{Data Sources and Inclusion Criteria}
% O*NET/task descriptors, labor stats, adoption proxies, education program attributes, etc.
\subsection{Data Preprocessing}
% missingness rules, normalization, feature construction, scaling.
\subsection{Constructed Indicators}
% e.g., Gen-AI exposure index, complementarity index, credential portability, etc.

\section{Model Overview}
\subsection{System Diagram and Modeling Logic}
% Fig.2 causal map: Gen-AI capability → task share shift → demand/supply → outcomes.
\subsection{Modular Architecture}
% (i) Career Future Module, (ii) Program Policy Module, (iii) Multi-metric Evaluation Module,
% (iv) Sensitivity/Robustness Module.

\section{Module I: Career-Future Projection Model (Three Careers)}
\subsection{Common Formulation (Reusable Across Careers)}
% Provide ONE unified formulation; then plug in career-specific parameters.
% Outputs: employment index, wage index, task mix, skill demand vector, uncertainty band.
\subsection{STEM Career Case}
\subsection{Trade Career Case}
\subsection{Arts Career Case}
% Each case: 3 paragraphs max:
% 1) conclusion, 2) method/parameters, 3) key quantitative results + figure reference.
% Figures: trajectory + uncertainty; task-mix shift; key drivers.

\section{Module II: Institution \& Program Response Optimization (Three Institutions)}
\subsection{Decision Levers}
% Enrollment size; curriculum hours; AI policy (ban/allow/teach); practicum structure; assessment redesign.
\subsection{Objective(s) and Constraints}
% Employability + optional success metrics (equity, creativity, sustainability, integrity).
\subsection{Solution Method}
% Multi-objective (Pareto / ε-constraint / weighted score) with transparent weights.
\subsection{Recommendations by Institution}
% University program (STEM), trade school program (Trade), arts school program (Arts).
% Deliver as an action table: action → mechanism → expected effect size → risk/mitigation.

\section{Module III: Beyond Employability (Optional, but high-scoring if concise)}
\subsection{Alternative Success Metrics}
% e.g., resilience, ethics/compliance, creative originality, societal value, sustainability footprint.
\subsection{How Recommendations Change Under Multi-metric Evaluation}
% Show trade-off plot (Pareto / spider chart).

\section{Generalization and Transferability}
% What generalizes across institutions/programs; boundary conditions; checklist.

\section{Sensitivity and Robustness}
\subsection{Local Sensitivity (SSC Tornado Ranking)}
% SSC = (ΔY/Y0)/(ΔX/X0)
\subsection{Interaction/Combination Effects}
% Heatmap or interaction curves: multi-parameter perturbations.
\subsection{Robustness Checks}
% Scenario stress tests / bootstrap / cross-validation (choose ≥1, visualize bands).

\section{Strengths and Weaknesses}
% 3 strengths + 3 weaknesses, each tied to judging criteria.

\section{Conclusions and Actionable Plan}
% 90-day / 1-year / 3-year roadmap per institution; link each action to a metric/figure.

\section*{References}
% Compact, high-quality, only what supports model/data choices.

\appendix
\section{Appendix (Optional, only if needed)}
% Derivations, extra tables, parameter calibration details.

% NOTE: If AI/LLM used:
\newpage
\section*{Report on Use of AI (Not counted in 25 pages)}
% What used for, model version, how verified, what not used for.
